% BHU PYQ -- Transcribed & Corrected
% Paper: PHIMN-31 / PHIMN31: Philosophical Counseling (Minor)
% Exam: B.A. (Hons.) Semester III Examination, 2025-26 (NEP)
% Category: Minor Core
% Department: Department of Philosophy, Faculty of Arts, Banaras Hindu University

\documentclass[11pt,a4paper]{article}
\usepackage[a4paper,top=2.5cm,bottom=2.5cm,left=2.8cm,right=2.8cm,headheight=14pt]{geometry}
\usepackage{amsmath,amssymb}
\usepackage[T1]{fontenc}
\usepackage[utf8]{inputenc}
\usepackage{lmodern,microtype,fancyhdr,enumitem,hyperref,lastpage,parskip}

\hypersetup{
    colorlinks=true,
    urlcolor=black,
    linkcolor=black,
    pdftitle={PHIMN31: Philosophical Counseling -- BHU B.A. Sem III 2025-26 (NEP)}
}

\pagestyle{fancy}
\fancyhf{}
\renewcommand{\headrulewidth}{0.4pt}
\renewcommand{\footrulewidth}{0.4pt}
\fancyhead[L]{\small \textit{Banaras Hindu University}}
\fancyhead[R]{\small \textit{PHIMN-31: Philosophical Counseling (Sem III 2025-26)}}
\fancyfoot[C]{\small Page \thepage\ of \pageref{LastPage}}

\newlist{parts}{enumerate}{2}
\setlist[parts,1]{label=(\alph*),leftmargin=*,itemsep=4pt,topsep=2pt}

\newcommand{\pts}[1]{\hfill \textit{[#1]}}

\begin{document}

\begin{center}
  {\large\bfseries BANARAS HINDU UNIVERSITY}\\[2pt]
  {\normalsize B.A. (Hons.) Semester III Examination, 2025-26 (NEP)}\\[6pt]
  \rule{0.8\linewidth}{0.5pt}\\[6pt]
  {\large\bfseries DEPARTMENT OF PHILOSOPHY}\\[4pt]
  {\large\bfseries Paper Code: PHIMN-31 : Philosophical Counseling (Minor)}\\[2pt]
  {\small Ref. Code: 2025-26/D-III/092}\\[4pt]
  {\normalsize Time: 2:30 Hours \hfill Full Marks: 70}
\end{center}

\rule{\linewidth}{0.4pt}

\smallskip

\noindent \textit{(Write your Roll No. at the top immediately on the receipt of this question paper)}\\[2pt]
\noindent \textit{Note: Attempt all the three Sections: A, B and C according to the instructions given.}\\[1pt]
\noindent \textit{दिये गये निर्देशानुसार सभी तीनों खण्डों: अ, ब तथा स के उत्तर दीजिए।}

\bigskip

\begin{center}
  {\large\bfseries SECTION A / खण्ड-अ}\\[2pt]
  {\normalsize (Long Answer Type Questions / दीर्घ उत्तरीय प्रश्न)}
\end{center}

\noindent \textit{Note: Attempt any three questions out of the following five questions, in about 300 words each. Each question carries 10 marks.}\\[2pt]
\noindent \textit{निम्नलिखित पाँच प्रश्नों में से किन्हीं तीन प्रश्नों के उत्तर दीजिए। प्रत्येक प्रश्न का उत्तर लगभग 300 शब्दों में दें। प्रत्येक प्रश्न 10 अंकों का है।}

\bigskip

\noindent \textbf{Question 1.} What is philosophical counseling? Discuss its nature and scope.\\[1pt]
दार्शनिक परामर्श क्या है? इसके स्वरूप एवं क्षेत्र की विवेचना कीजिए। \pts{10}

\medskip\rule{\linewidth}{0.2pt}\medskip

\noindent \textbf{Question 2.} Discuss the difference between psychological counseling and philosophical counseling.\\[1pt]
मनोवैज्ञानिक परामर्श तथा दार्शनिक परामर्श के मध्य अंतर की विवेचना कीजिए। \pts{10}

\medskip\rule{\linewidth}{0.2pt}\medskip

\noindent \textbf{Question 3.} What is anxiety? Can philosophical counseling play a role in reducing it? Discuss.\\[1pt]
चिंता क्या है? क्या चिंता को कम करने में दार्शनिक परामर्श की भूमिका हो सकती है? विवेचना कीजिए। \pts{10}

\medskip\rule{\linewidth}{0.2pt}\medskip

\noindent \textbf{Question 4.} What is the role of philosophy in critical thinking? Explain with suitable examples.\\[1pt]
आलोचनात्मक चिंतन में दर्शन की क्या भूमिका है? उपयुक्त उदाहरणों के साथ स्पष्ट कीजिए। \pts{10}

\medskip\rule{\linewidth}{0.2pt}\medskip

\noindent \textbf{Question 5.} What is an ethical dilemma, and how can philosophical counseling help in resolving it? Discuss.\\[1pt]
नैतिक दुविधा क्या है, तथा दार्शनिक परामर्श इसे सुलझाने में किस प्रकार सहायता कर सकता है? विवेचना कीजिए। \pts{10}

\bigskip
\rule{\linewidth}{0.2pt}
\bigskip

\begin{center}
  {\large\bfseries SECTION B / खण्ड-ब}\\[2pt]
  {\normalsize (Short Answer Type Questions / लघु उत्तरीय प्रश्न)}
\end{center}

\noindent \textit{Note: Attempt any four questions out of the following six questions, in about 150 words each. Each question carries 5 marks.}\\[2pt]
\noindent \textit{निम्नलिखित छः प्रश्नों में से किन्हीं चार प्रश्नों के उत्तर दीजिए। प्रत्येक प्रश्न का उत्तर लगभग 150 शब्दों में दें। प्रत्येक प्रश्न 5 अंकों का है।}

\bigskip

\noindent \textbf{Question 6.} How does philosophical counseling differ from spiritual counseling?\\[1pt]
दार्शनिक परामर्श, आध्यात्मिक परामर्श से किस प्रकार भिन्न है? \pts{5}

\medskip\rule{\linewidth}{0.2pt}\medskip

\noindent \textbf{Question 7.} Write a short note on the concept of ``Socratic Dialogue'' in counseling.\\[1pt]
परामर्श में ``सुकराती संवाद'' की अवधारणा पर संक्षिप्त टिप्पणी लिखिए। \pts{5}

\medskip\rule{\linewidth}{0.2pt}\medskip

\noindent \textbf{Question 8.} Discuss the role of Stoic philosophy in coping with life challenges.\\[1pt]
जीवन की चुनौतियों से निपटने में स्टोइक दर्शन की भूमिका की विवेचना कीजिए। \pts{5}

\medskip\rule{\linewidth}{0.2pt}\medskip

\noindent \textbf{Question 9.} What is the significance of self-examination in philosophical counseling?\\[1pt]
दार्शनिक परामर्श में आत्म-परीक्षण का क्या महत्व है? \pts{5}

\medskip\rule{\linewidth}{0.2pt}\medskip

\noindent \textbf{Question 10.} Explain the concept of ``Meaningfulness of Life'' in existential counseling.\\[1pt]
अस्तित्ववादी परामर्श में ``जीवन की अर्थपूर्णता'' की अवधारणा को स्पष्ट कीजिए। \pts{5}

\medskip\rule{\linewidth}{0.2pt}\medskip

\noindent \textbf{Question 11.} Write a short note on PEACE methodology in philosophical counseling.\\[1pt]
दार्शनिक परामर्श में PEACE कार्यप्रणाली पर एक संक्षिप्त टिप्पणी लिखिए। \pts{5}

\bigskip
\rule{\linewidth}{0.2pt}
\bigskip

\begin{center}
  {\large\bfseries SECTION C / खण्ड-स}\\[2pt]
  {\normalsize (Very Short Answer Type Questions / अतिलघु उत्तरीय प्रश्न)}
\end{center}

\noindent \textit{Note: Answer all ten parts of the following Question No. 12 in not more than two sentences. Each part carries 2 marks.}\\[2pt]
\noindent \textit{अधोलिखित प्रश्न संख्या 12 के सभी दस भागों के उत्तर अधिकतम दो वाक्यों में दीजिए। प्रत्येक भाग 2 अंकों का है।}

\bigskip

\noindent \textbf{Question 12.}
\begin{parts}
  \item What is the main method used in philosophical counseling?\\[1pt]
        दार्शनिक परामर्श में प्रयोग की जाने वाली प्रमुख विधि क्या है? \pts{2}

  \item Who is called the Father of Existentialism?\\[1pt]
        अस्तित्ववाद का जनक किसे कहा जाता है? \pts{2}

  \item For whom is philosophical counseling useful?\\[1pt]
        दार्शनिक परामर्श किसके लिए उपयोगी है? \pts{2}

  \item What is the full form of ``APPA''?\\[1pt]
        ``APPA'' का पूर्ण रूप क्या है? \pts{2}

  \item What is the main tool of a philosophical counsellor?\\[1pt]
        दार्शनिक परामर्शदाता का मुख्य उपकरण क्या है? \pts{2}

  \item Who said ``Live according to nature''?\\[1pt]
        ``प्रकृति के अनुसार जियो'' यह किसने कहा था? \pts{2}

  \item Who said that ``man is condemned to be free''?\\[1pt]
        किसने कहा है कि ``मनुष्य स्वतंत्र होने के लिए अभिशप्त है''? \pts{2}

  \item What is the full form of ``CBT''?\\[1pt]
        ``सीबीटी'' का पूर्ण रूप क्या है? \pts{2}

  \item Who initiated philosophical counseling?\\[1pt]
        दार्शनिक परामर्श की शुरूआत किसने की? \pts{2}

  \item Name any one founder of modern philosophical counseling.\\[1pt]
        आधुनिक दार्शनिक परामर्श के किसी एक संस्थापक का नाम लिखिए। \pts{2}
\end{parts}

\vfill
\rule{\linewidth}{0.4pt}

\begin{center}\small
\href{https://bhu-exam.vercel.app/}{Click here to visit our website}\\[4pt]
\href{https://me-aryan.vercel.app/}{Click here to visit our contributor}
\end{center}

\end{document}
